\section{Introduction}

Systematic reviews have become a dominant method and, in many areas, the standard for synthesizing evidence to answer focused research questions, especially in medicine, education, and the social sciences. Systematic reviews follow a protocolled and transparent approach to searching for, selecting, and documenting relevant literature, with explicit inclusion and exclusion criteria. While they share this systematic approach to evidence identification and screening, they can differ substantially with respect to synthesis methods, ranging from qualitative evidence syntheses and mixed-methods reviews to quantitative meta-analyses that aggregate effect sizes~\cite{gough:2012}. Systematic reviews are valued for their completeness in capturing the available evidence, transparency in documenting procedures and data sources, reproducibility through well-defined methodologies, rigorous appraisal of study quality (including risk-of-bias assessment and related threats to validity), and their ability to support more robust and informative conclusions via structured synthesis across studies~\cite{liberati:2009, mulrow:1994}. Their results guide subsequent research, policy, and professional practice.

Historically, systematic review methodology was developed and institutionalized most prominently in biomedicine to extract and communicate ``the evidence'' from an expanding primary literature for decision-making~\cite{chalmers:1993, guyatt:1992}. Related evidence-based policy and practice agendas then helped diffuse evidence-synthesis toolkits beyond health, supported by new institutions and infrastructures for producing and disseminating reviews~\cite{simons:2021}. Over time, these principles spread into many other disciplines, often with field-specific adaptations in information sources, search conventions, eligibility criteria, and synthesis approaches. Regardless of the domain and application, a systematic review is typically conducted in a multi-step workflow, from scoping the objective and inclusion criteria, through searching for and selecting studies (retrieval and screening), collecting and appraising study data, and finally synthesizing, presenting, and interpreting findings. Methodological guides in biomedicine~\cite{chandler_cochrane_2019}, nursing~\cite{bettany_saltikov_how_2012}, education~\cite{newman_systematic_2020}, social science~\cite{petticrew_systematic_2008}, design~\cite{lame:2019}, or computer science~\cite{kitchenham:2004} cover these core steps, while differing in granularity and how they group them into phases.

Conducting a systematic review is time-consuming and expensive. Consequently, substantial effort has been devoted to automating one of the most expensive parts of the process, namely retrieval and screening. However, existing evaluation resources for retrieval and screening remain largely concentrated in the biomedical domain (see Table~\ref{table-sr-datasets}). This focus reflects the empirical dominance of systematic reviews in biomedicine, but constrains evaluation in other scientific fields. Consequently, domain-specific query formulations, search strategies, and relevance criteria outside biomedicine remain sparsely represented in current collections, limiting the assessment of retrieval and screening methods beyond the medical domain.

\newcommand{\bsq}{$\blacksquare$}
\newcommand{\wsq}{$\square$}


\begin{table*}[t]
\small
\sffamily
\renewcommand{\arraystretch}{1.05}
\renewcommand{\tabcolsep}{3.6pt}
\centering
\caption{Overview of systematic review datasets.}
\label{table-sr-datasets}
\vskip-2ex
\begin{tabular}{@{}l@{}rcc@{\hspace{2em}}lcccc@{\hspace{2em}}rcccccc@{\hspace{2em}}ccccc@{}}
\toprule
  \multicolumn{4}{@{}c@{\hspace{2em}}}{\bfseries Dataset}                                                                                                   & \multicolumn{5}{@{}c@{\hspace{2em}}}{\bfseries Publication Data} & \multicolumn{12}{@{}c@{}}{\bfseries Structured Data Fields}    \\
\cmidrule(r@{2em}){1-4}\cmidrule(r@{2em}){5-9}\cmidrule{10-21}
  Name             & \multicolumn{1}{@{}c@{}}{Rf.} & Year               & Link                                         & Domain & \multicolumn{3}{@{}c@{}}{\kern-0.5em SysRevs} & Pool & Topics & \multicolumn{6}{@{}c@{\hspace{2em}}}{Topic Description} & \multicolumn{5}{@{}c@{}}{Screening} \\
\cmidrule(r@{\tabcolsep}){1-1}\cmidrule(r@{\tabcolsep}){2-2}\cmidrule(l@{\tabcolsep}r@{\tabcolsep}){3-3}\cmidrule(l@{\tabcolsep}r@{2em}){4-4}\cmidrule(r@{\tabcolsep}){5-5}\cmidrule(l@{\tabcolsep}r@{\tabcolsep}){6-8}\cmidrule(l@{\tabcolsep}r@{2em}){9-9}\cmidrule(r@{\tabcolsep}){10-10}\cmidrule(l@{\tabcolsep}r@{2em}){11-16}\cmidrule{17-21}
                   &                        &                           &                                                                                       &         & Rfs. &  TA  &  FT  &          &        & Obj. & RQs  & KWs  & BQs  & Crt. & Rng. & DBs  &  CR  &  CI  &  PR  &  PI  \\
\midrule
  Cohen et al.     & \cite{cohen:2006}      & \citeyear{cohen:2006}     & \href{https://dmice.ohsu.edu/cohenaa/systematic-drug-class-review-data.html}{\faLink} & Biomed  &      &      &      &          &     15 &      &      &      &      &      &      &      & \bsq & \bsq & \bsq & \bsq \\
  Wallace et al.   & \cite{wallace:2010}    & \citeyear{wallace:2010}   & \href{https://github.com/bwallace/citation-screening}{\faLink}                        & Biomed  &   \bsq   &      &      &          &      3 &      &      &      &      &      &      &      & \bsq & \bsq &      &      \\
  SWIFT-Review     & \cite{howard:2016}     & \citeyear{howard:2016}    & \href{https://link.springer.com/article/10.1186/s13643-016-0263-z#Sec30}{\faLink}     & Biomed  &      &      &      &          &      5 &      &      &      &      &      &      &      & \bsq & \bsq & \bsq & \bsq \\
  SysRev Query C.  & \ \cite{scells:2017}   & \citeyear{scells:2017}    & \href{https://github.com/ielab/SIGIR2017-SysRev-Collection}{\faLink}                  & Biomed  &   \bsq   &      &      &     PubMed     &     94 &      &      &      & \bsq &      & \bsq &      & \bsq & \bsq & \bsq & \bsq \\
  CLEF TAR 2017    & \cite{kanoulas:2017}   & \citeyear{kanoulas:2017}  & \href{https://github.com/CLEF-TAR/tar}{\faLink}                                       & Biomed  &   \bsq    &  \wsq    &      &      PubMed    &     50 &      &      &      & \bsq &      & \bsq &      & \bsq & \bsq & \bsq & \bsq \\
  CLEF TAR 2018    & \cite{kanoulas:2018}   & \citeyear{kanoulas:2018}  & \href{https://github.com/CLEF-TAR/tar}{\faLink}                                       & Biomed  &  \bsq     &   \wsq   &      &      PubMed    &     30 &      &      &      & \bsq &      & \bsq &      & \bsq & \bsq & \bsq & \bsq \\
  CLEF TAR 2019    & \cite{kanoulas:2019}   & \citeyear{kanoulas:2019}  & \href{https://github.com/CLEF-TAR/tar}{\faLink}                                       & Biomed  &   \bsq    &   \wsq   &      &     PubMed     &     49 &      &      &      & \bsq &      & \bsq &      & \bsq & \bsq & \bsq & \bsq \\
  Alharbi et al.   & \cite{alharbi:2019}    & \citeyear{alharbi:2019}   & \href{https://github.com/Amal-Alharbi/Systematic_Reviews_Update}{\faLink}             & Biomed  &  \bsq   &  \wsq    &      &      PubMed   &     25 &      &      &   \bsq   &   \bsq   &      &      &      &      &   \bsq   &      &   \bsq   \\
  Hannousse et al. & \cite{hannousse:2021}  & \citeyear{hannousse:2021} & \href{https://github.com/hannousse/Semantic-Scholar-Evaluation}{\faLink}              & CompSci &  \bsq    &      &      &     S2     &      7 &      &      &      &      &      &      &      &      &  \bsq    &      &  \bsq    \\
  SysRev Seed C.   & \cite{wang:2022}       & \citeyear{wang:2022}      & \href{https://github.com/ielab/sysrev-seed-collection}{\faLink}                       & Biomed  &   \bsq   &  \wsq    &      &     PubMed     &     40 & \bsq &      &      & \bsq &      & \bsq &      & \bsq & \bsq & \bsq & \bsq \\
  CSMeD            & \cite{kusa:2023}       & \citeyear{kusa:2023}      &  \href{https://github.com/WojciechKusa/systematic-review-datasets}{\faLink}            & Both    &   \bsq   &  \wsq    &      &   Mixed  &    325 & \bsq &      &      & \bsq & \bsq &      & \bsq & \bsq & \bsq & \bsq & \bsq \\
  AutoBool         & \cite{wang:2025b}      & \citeyear{wang:2025b}     & \href{https://github.com/ielab/AutoBool}{\faLink}                                     & Biomed  &   \bsq   &  \wsq    &      &     PubMed     & 65,588 &      &      &      &      &      & \bsq &      &      & \bsq &      & \bsq \\
\midrule
  \multicolumn{2}{@{}l@{}}{Webis-SR4ALL-26} & 2026                      & \href{https://doi.org/10.5281/zenodo.18431943}{\faLink}                               & Open    & \bsq & \bsq & \wsq & OpenAlex & 301,871 & \wsq & \wsq & \wsq & \wsq & \wsq & \wsq & \wsq & \wsq & \wsq &      & \bsq \\
\bottomrule
\addlinespace
\multicolumn{21}{@{}l@{}}{
\parbox{\textwidth}{
\footnotesize Abbreviations: For systematic reviews (SysRevs), we distinguish whether they are available as references (Rfs.), title and abstract (TA), and full texts (FT). For the structured data fields, we distinguish objective~(Obj.), research questions~(RQs), keywords~(KWs), Boolean queries~(BQs), inclusion and exclusion criteria~(Crt.), publication date ranges~(Rng.), publication databases originally queried~(DBs), counts of retrieved~(CR) and included~(CI) publications, and references to publications retrieved~(PR) and included~(PI) in the systematic review. Filled boxes indicate full availability; empty boxes indicate partial availability.}}
\end{tabular}
\end{table*}


To address this gap, we introduce Webis-SR4ALL-26, a large-scale, cross-disciplinary corpus of 301,871~systematic reviews spanning 27~scientific domains, derived from reviews indexed in OpenAlex~\cite{Priem2022OpenAlexAF}. A large subset of reviews is linked to its resolved set of cited studies, providing a stable bibliographic context. In addition, a broad range of methodological information is extracted from the available review full texts when explicitly reported. For systematic reviews that report a search strategy, whether as explicit Boolean queries or keyword lists, we construct a canonical Boolean representation as a minimal approximation of the reported strategy for execution in OpenAlex. Together, these components make Webis-SR4ALL-26 a rich multi-layer resource for retrieval and screening in systematic reviews across domains and review types, and for multiple research communities. Its layers (links, method fields, normalized queries) can each serve as supervision, evaluation targets, or observational signals for descriptive analyses. The corpus supports benchmarking and development of retrieval, screening, and extraction approaches, and enables comparative meta-research, including science-of-science and sociological studies of how query formulation and reporting practices evolve across domains and review styles over time.
% The collection is available under the Creative Commons Attribution 4.0 International (CC~BY~4.0) license, and accompanying code is released under the MIT~License to support reproducible research.




%Retrieval is often operationalized by translating the review objective into a complex Boolean query, typically guided by one or more research questions~\cite{scells2021comparison}. Boolean search supports interpretability (i.e. the Boolean retrieval model is straightforward to understand and reason about), and when combined with appropriate temporal restrictions and stable databases (i.e. the stored documents remain unchanged), enables reproducible reruns of the search strategy~\cite{staudinger2025timir}. These queries define the pool of eligible candidate studies and directly constrain the outcome of the review, as recall losses incurred at retrieval are difficult to recover in later stages~\cite{MACFARLANE2022200091}. Screening then selects the eligible studies from this candidate set by applying the review's eligibility criteria, often under strong recall requirements, and it constitutes a major share of the human effort in producing a review. Because retrieval errors are difficult to recover downstream and screening effort is substantial, there is a need for large-scale, cross-domain resources that support reproducible evaluation of retrieval and screening methods.

%Prior work on automatic Boolean query generation ranges from computational adaptations of expert-designed strategies~\cite{scells2021comparison} to more recent investigations using large language models (LLMs)~\cite{wang2023can,wang2025reassessing}. While LLM-based approaches show promise, they continue to involve trade-offs in precision and recall that are sensitive to evaluation settings~\cite{DBLP:conf/sigir/WangSKZ23}. Recent work has further shown that Boolean query generation can be optimized end-to-end using reinforcement learning by directly rewarding retrieval effectiveness, highlighting both the feasibility of learned high-recall strategies and the growing need for large-scale, executable evaluation resources \cite{autobool}. LLMs have also been applied to the screening task, where they achieve promising zero-shot performance when explicitly calibrated for recall~\cite{DBLP:conf/ecir/WangSZPKZ24}. Beyond individual components, agent-based systems have been proposed to coordinate multiple stages of the systematic review workflow, further increasing the reliance on realistic and reproducible evaluation resources~\cite{DBLP:journals/corr/abs-2403-08399}.





